% !TEX root = ../metatime_monograph.tex
\chapter{Poincar\'{e} Disk Geometry}
\label{app:poincare}

% CITATION_CONTEXT_V10_9_BEGIN
\paragraph{Established context.}
The differential-geometric and topological language used for disk geometry,
connections, curvature, and bundles follows standard references
\cite{KobayashiNomizu1963,BottTu1982,Nakahara2003}.  The embedding of the
Metatime state construction into this geometry is specific to the present model.
% CITATION_CONTEXT_V10_9_END

\section{Definition}

The Poincar\'{e} disk $\mathbb{D}$ is the unit disk in the complex plane
equipped with the metric:

\[
ds^2 = \frac{4\,|dz|^2}{(1 - |z|^2)^2},
\]

where $z \in \mathbb{C}$, $|z| < 1$.  This is a model of the
two-dimensional hyperbolic plane with constant negative curvature
$K = -1$.

\section{Isometries}

The isometry group of $\mathbb{D}$ is $\text{PSU}(1,1)$, consisting of
M\"{o}bius transformations:

\[
z \mapsto e^{i\theta}\,\frac{z - a}{1 - \bar{a}z},
\qquad a \in \mathbb{D},\ \theta \in [0,2\pi).
\]

\section{Geodesics}

Geodesics in $\mathbb{D}$ are arcs of circles orthogonal to the boundary
$|z| = 1$, or diameters through the origin.  The geodesic distance between
two points $z_1$ and $z_2$ is:

\[
d(z_1, z_2) = \operatorname{artanh}\left|
\frac{z_1 - z_2}{1 - \bar{z}_1 z_2}
\right|.
\]

\section{Boundary and Compactification}

The boundary $\partial\mathbb{D} = \{z : |z| = 1\}$ corresponds to the
conformal infinity of the hyperbolic plane.  In the Metatime framework,
the assignment of information degrees of freedom to this boundary is a model
interpretation rather than a theorem of hyperbolic geometry.

\section{Kuramoto Model on $\mathbb{D}$}

A model system of $N$ coupled oscillators on $\mathbb{D}$ is written here as

\[
\dot{z}_i = (1 - |z_i|^2)\left(i\omega_i z_i + \sum_{j\neq i} K_{ij} z_j\right),
\]

where $\omega_i$ are natural frequencies and $K_{ij}$ are model coupling
strengths.  Specific assignments of those couplings to quark-prime data belong
to the Metatime ansatz, not to the standard Kuramoto model itself.

The synchronization order parameter is:

\[
r e^{i\psi} = \frac{1}{N}\sum_{j=1}^N \frac{z_j}{|z_j|}.
\]

For mesonic applications the framework studies $N=2$ sectors and relates the
resulting trajectories to action coordinates.  Any mass-gap interpretation
remains part of the model layer and is tested elsewhere in the monograph.

\section{Hyperbolic area and the $\kappa$ obligation}

For curvature $K=-1$, the area of a geodesic hyperbolic triangle with interior
angles $\alpha,\beta,\gamma$ is

\[
A = \pi - (\alpha + \beta + \gamma).
\]

Earlier drafts stated that a tetrahedral angle involving
\[
\alpha = \arccos\frac{L_4}{L_3} = \arccos\frac{2}{7}
\]
``gives $\kappa/2$''.  That statement is incomplete: a triangle area is not
determined by one angle, and no theorem follows until the remaining angles,
geodesic construction, curvature normalization, and equality to $\kappa/2$
are all specified independently.

Accordingly, the publication-level statement is now the explicit obligation
\[
\boxed{
\text{derive a complete hyperbolic cell }(\alpha,\beta,\gamma)
\text{ and prove }A=\kappa/2
}
\]
if such a bridge is intended.  Until that construction is supplied, the
Poincar\'e geometry and the Metatime value of $\kappa$ are cross-referenced but
not identified by this area formula.  The current normalization status is
summarized in Appendix~\ref{app:kappa} and the paired dependency map in
Appendix~\ref{app:information_spinor_crosswalk}.
